% GENERATED FILE. Edit canonical lessons, then run npm run generate:books.
% canonical-source-hash: fnv1a64:7adc390977f5b694
% canonical-lessons: SA-C339-yuvarajah, SA-C339-svayamvarah, SA-C339-vrnite, SA-C339-pracchannah, SA-C339-nyayah, SA-R339-review-333-339, SA-R339-review-final-339

\chapter{The Prince's Choice}
\label{ch:sa-the-prince-s-choice}
\hlchaptermodality{339}

\begin{chapteropening}
\textbf{By the end of this chapter:} \emph{I can tell how a crown prince won a bride.}

Complete the last lesson of chapter 339: I can tell how a crown prince won a bride.
\end{chapteropening}

\section[\texorpdfstring{\sk{युवराजः} (yuvarājaḥ)}{युवराजः (yuvarājaḥ)}]{\sk{युवराजः} (yuvarājaḥ) — a crown prince}

\label{lesson:SA-C339-yuvarajah}



\begin{quote}
Before the new one: say the Sanskrit for hell, then the Sanskrit for reaps.
\end{quote}

\subsection*{You'll want to know: \sk{युवराजः}}
\textbf{\sk{युवराजः}} — \emph{yuvarājaḥ} — \textquotedblleft{}a crown prince\textquotedblright{}.

\begin{quote}
\sk{युवराजः} \sk{अश्वेन} \sk{नगरम्} \sk{आगच्छति।}
\end{quote}

\begin{quote}
\emph{yuvarājaḥ aśvena nagaram āgacchati.}
\end{quote}

\begin{quote}
The crown prince comes to the city on horseback.
\end{quote}

A masculine noun.

\begin{cousinweb}
\textbf{\sk{युवराजः}} (\emph{yuvarājaḥ}) is built from \emph{yuvan}, \textquotedblleft{}young\textquotedblright{}, and \emph{rājan}, \textquotedblleft{}a king\textquotedblright{}; MW gives \emph{yuvarāja} as a young king, an heir-apparent, a crown prince (MW).
\end{cousinweb}

\begin{grammarlens}[title={What you've built}]
One more word for telling what happened.
\end{grammarlens}

\subsection*{Your turn}
\begin{itemize}
\raggedright
  \item \emph{Say it:} \emph{yuvarājaḥ}
  \item \emph{Say it:} \emph{yuvarājaḥ}, once more
  \item \emph{Say it:} \emph{yuvarājaḥ aśvena nagaram āgacchati.} — The crown prince comes to the city on horseback.
  \item \emph{Recall:} say the Sanskrit for hell, then the Sanskrit for reaps, then say \emph{yuvarājaḥ} again
\end{itemize}

\begin{tcolorbox}[breakable,colback=teal!4,colframe=teal!35!black,title={Before you move on}]
What does \sk{युवराजः} mean? (\textquotedblleft{}A crown prince\textquotedblright{}.) Say it once more.
\end{tcolorbox}
\section[\texorpdfstring{\sk{स्वयंवरः} (svayaṃvaraḥ)}{स्वयंवरः (svayaṃvaraḥ)}]{\sk{स्वयंवरः} (svayaṃvaraḥ) — a bride's choice of a husband}

\label{lesson:SA-C339-svayamvarah}



\begin{quote}
Before the new one: say the Sanskrit for reaps, then the Sanskrit for a crown prince.
\end{quote}

\subsection*{You'll want to know: \sk{स्वयंवरः}}
\textbf{\sk{स्वयंवरः}} — \emph{svayaṃvaraḥ} — \textquotedblleft{}a bride's choice of a husband\textquotedblright{}.

\begin{quote}
\sk{श्वः} \sk{प्रासादे} \sk{स्वयंवरः} \sk{भविष्यति।}
\end{quote}

\begin{quote}
\emph{śvaḥ prāsāde svayaṃvaraḥ bhaviṣyati.}
\end{quote}

\begin{quote}
Tomorrow a princess will choose her husband at the palace.
\end{quote}

A masculine noun.

\begin{cousinweb}
\textbf{\sk{स्वयंवरः}} (\emph{svayaṃvaraḥ}) is built from \emph{svayam}, \textquotedblleft{}self\textquotedblright{}, and \emph{vara}, \textquotedblleft{}choosing\textquotedblright{}, from the root √vṛ; MW gives \emph{svayaṃvara} as self-choice, the choosing of a husband by a princess at an assembly of suitors (MW).
\end{cousinweb}

\begin{grammarlens}[title={What you've built}]
One more word for telling what happened.
\end{grammarlens}

\subsection*{Your turn}
\begin{itemize}
\raggedright
  \item \emph{Say it:} \emph{svayaṃvaraḥ}
  \item \emph{Say it:} \emph{svayaṃvaraḥ}, once more
  \item \emph{Say it:} \emph{śvaḥ prāsāde svayaṃvaraḥ bhaviṣyati.} — Tomorrow a princess will choose her husband at the palace.
  \item \emph{Recall:} say the Sanskrit for reaps, then the Sanskrit for a crown prince, then say \emph{svayaṃvaraḥ} again
\end{itemize}

\begin{tcolorbox}[breakable,colback=teal!4,colframe=teal!35!black,title={Before you move on}]
What does \sk{स्वयंवरः} mean? (\textquotedblleft{}A bride's choice of a husband\textquotedblright{}.) Say it once more.
\end{tcolorbox}
\section[\texorpdfstring{\sk{वृणीते} (vṛṇīte)}{वृणीते (vṛṇīte)}]{\sk{वृणीते} (vṛṇīte) — chooses}

\label{lesson:SA-C339-vrnite}



\begin{quote}
Before the new one: say the Sanskrit for a crown prince, then the Sanskrit for a bride's choice.
\end{quote}

\subsection*{You'll want to know: \sk{वृणीते}}
\textbf{\sk{वृणीते}} — \emph{vṛṇīte} — \textquotedblleft{}chooses\textquotedblright{}.

\begin{quote}
\sk{कुमारी} \sk{युवराजं} \sk{वृणीते।}
\end{quote}

\begin{quote}
\emph{kumārī yuvarājaṃ vṛṇīte.}
\end{quote}

\begin{quote}
The maiden chooses the crown prince.
\end{quote}

Past (she …): \sk{सा} \sk{युवराजं} \sk{वृतवती}, \emph{sā yuvarājaṃ vṛtavatī}.

\begin{cousinweb}
\textbf{\sk{वृणीते}} (\emph{vṛṇīte}) comes from the root √vṛ, \textquotedblleft{}to choose\textquotedblright{} (MW's second √vṛ), the root of \emph{vara} too; MW gives \emph{vṛta} as chosen, selected (MW).
\end{cousinweb}

\begin{grammarlens}[title={What you've built}]
One more word for telling what happened.
\end{grammarlens}

\subsection*{Your turn}
\begin{itemize}
\raggedright
  \item \emph{Say it:} \emph{vṛṇīte}
  \item \emph{Say it:} \emph{vṛṇīte}, once more
  \item \emph{Say it:} \emph{kumārī yuvarājaṃ vṛṇīte.} — The maiden chooses the crown prince.
  \item \emph{Recall:} say the Sanskrit for a crown prince, then the Sanskrit for a bride's choice, then say \emph{vṛṇīte} again
\end{itemize}

\begin{tcolorbox}[breakable,colback=teal!4,colframe=teal!35!black,title={Before you move on}]
What does \sk{वृणीते} mean? (\textquotedblleft{}Chooses\textquotedblright{}.) Say it once more.
\end{tcolorbox}
\section[\texorpdfstring{\sk{प्रच्छन्नः} (pracchannaḥ)}{प्रच्छन्नः (pracchannaḥ)}]{\sk{प्रच्छन्नः} (pracchannaḥ) — disguised; hidden}

\label{lesson:SA-C339-pracchannah}



\begin{quote}
Before the new one: say the Sanskrit for a bride's choice, then the Sanskrit for chooses.
\end{quote}

\subsection*{You'll want to know: \sk{प्रच्छन्नः}}
\textbf{\sk{प्रच्छन्नः}} — \emph{pracchannaḥ} — \textquotedblleft{}disguised; hidden\textquotedblright{}.

\begin{quote}
\sk{युवराजः} \sk{प्रच्छन्नः} \sk{नगरे} \sk{भ्रमति।}
\end{quote}

\begin{quote}
\emph{yuvarājaḥ pracchannaḥ nagare bhramati.}
\end{quote}

\begin{quote}
The crown prince wanders the city in disguise.
\end{quote}

\begin{cousinweb}
\textbf{\sk{प्रच्छन्नः}} (\emph{pracchannaḥ}) is built from \emph{pra} and \emph{channa}, \textquotedblleft{}covered\textquotedblright{}, from the root √chad, \textquotedblleft{}to cover\textquotedblright{}; MW gives \emph{pracchanna} as covered, hidden, secret, disguised (MW).
\end{cousinweb}

\begin{grammarlens}[title={What you've built}]
One more word for telling what happened.
\end{grammarlens}

\subsection*{Your turn}
\begin{itemize}
\raggedright
  \item \emph{Say it:} \emph{pracchannaḥ}
  \item \emph{Say it:} \emph{pracchannaḥ}, once more
  \item \emph{Say it:} \emph{yuvarājaḥ pracchannaḥ nagare bhramati.} — The crown prince wanders the city in disguise.
  \item \emph{Recall:} say the Sanskrit for a bride's choice, then the Sanskrit for chooses, then say \emph{pracchannaḥ} again
\end{itemize}

\begin{tcolorbox}[breakable,colback=teal!4,colframe=teal!35!black,title={Before you move on}]
What does \sk{प्रच्छन्नः} mean? (\textquotedblleft{}Disguised; hidden\textquotedblright{}.) Say it once more.
\end{tcolorbox}
\section[\texorpdfstring{\sk{न्यायः} (nyāyaḥ)}{न्यायः (nyāyaḥ)}]{\sk{न्यायः} (nyāyaḥ) — justice; fairness}

\label{lesson:SA-C339-nyayah}



\begin{quote}
Before the new one: say the Sanskrit for chooses, then the Sanskrit for disguised.
\end{quote}

\subsection*{You'll want to know: \sk{न्यायः}}
\textbf{\sk{न्यायः}} — \emph{nyāyaḥ} — \textquotedblleft{}justice; fairness\textquotedblright{}.

\begin{quote}
\sk{युवराजः} \sk{न्यायेन} \sk{राज्यं} \sk{शास्ति।}
\end{quote}

\begin{quote}
\emph{yuvarājaḥ nyāyena rājyaṃ śāsti.}
\end{quote}

\begin{quote}
The crown prince rules the kingdom with justice.
\end{quote}

A masculine noun.

\begin{cousinweb}
\textbf{\sk{न्यायः}} (\emph{nyāyaḥ}) is built from \emph{ni} and the root √i, \textquotedblleft{}to go\textquotedblright{}: MW gives \emph{nyāya} as a rule or method, a right or fit manner, propriety, and a lawsuit or judicial sentence, and justice is its everyday modern sense (MW).
\end{cousinweb}

\begin{grammarlens}[title={What you've built}]
One more word for telling what happened.
\end{grammarlens}

\subsection*{Your turn}
\begin{itemize}
\raggedright
  \item \emph{Say it:} \emph{nyāyaḥ}
  \item \emph{Say it:} \emph{nyāyaḥ}, once more
  \item \emph{Say it:} \emph{yuvarājaḥ nyāyena rājyaṃ śāsti.} — The crown prince rules the kingdom with justice.
  \item \emph{Recall:} say the Sanskrit for chooses, then the Sanskrit for disguised, then say \emph{nyāyaḥ} again
\end{itemize}

\begin{tcolorbox}[breakable,colback=teal!4,colframe=teal!35!black,title={Before you move on}]
What does \sk{न्यायः} mean? (\textquotedblleft{}Justice; fairness\textquotedblright{}.) Say it once more.
\end{tcolorbox}
\section[(dialogue)]{Looking back: from The Curse and the Boon to The Prince's Choice}

\label{lesson:SA-R339-review-333-339}



\begin{quote}
Say the Sanskrit for a curse and a god.
\end{quote}

\subsection*{Your turn}
Say these aloud:

\begin{itemize}
\raggedright
  \item a curse, a god, a boon, shows favour, immortal
  \item a demon, an enemy of the gods, the slaying, kills, swallows
  \item illusion, a disguise, a trick, bewilders, attracts
  \item wise conduct, greed, a sea monster, bewildered, tempts
  \item life in the forest, a promise, a sandal, promises, hunting
  \item merit, a sin, heaven, hell, reaps
  \item a crown prince, a bride's choice, chooses, disguised, justice
\end{itemize}

\begin{tcolorbox}[breakable,colback=teal!4,colframe=teal!35!black,title={Before you move on}]
A curse? (\textbf{\sk{शापः}}, \emph{śāpaḥ}.) A god? (\textbf{\sk{देवः}}, \emph{devaḥ}.) A boon? (\textbf{\sk{वरः}}, \emph{varaḥ}.) Shows favour? (\textbf{\sk{प्रसीदति}}, \emph{prasīdati}.) Immortal? (\textbf{\sk{अमरः}}, \emph{amaraḥ}.) A demon? (\textbf{\sk{राक्षसः}}, \emph{rākṣasaḥ}.) An enemy of the gods? (\textbf{\sk{असुरः}}, \emph{asuraḥ}.) The slaying? (\textbf{\sk{वधः}}, \emph{vadhaḥ}.) Kills? (\textbf{\sk{हन्ति}}, \emph{hanti}.) Swallows? (\textbf{\sk{निगिरति}}, \emph{nigirati}.) Illusion? (\textbf{\sk{माया}}, \emph{māyā}.) A disguise? (\textbf{\sk{वेषः}}, \emph{veṣaḥ}.) A trick? (\textbf{\sk{छलम्}}, \emph{chalam}.) Bewilders? (\textbf{\sk{मोहयति}}, \emph{mohayati}.) Attracts? (\textbf{\sk{आकर्षति}}, \emph{ākarṣati}.) Wise conduct? (\textbf{\sk{नीतिः}}, \emph{nītiḥ}.) Greed? (\textbf{\sk{लोभः}}, \emph{lobhaḥ}.) A sea monster? (\textbf{\sk{मकरः}}, \emph{makaraḥ}.) Bewildered? (\textbf{\sk{मूढः}}, \emph{mūḍhaḥ}.) Tempts? (\textbf{\sk{लोभयति}}, \emph{lobhayati}.) Life in the forest? (\textbf{\sk{वनवासः}}, \emph{vanavāsaḥ}.) A promise? (\textbf{\sk{प्रतिज्ञा}}, \emph{pratijñā}.) A sandal? (\textbf{\sk{पादुका}}, \emph{pādukā}.) Promises? (\textbf{\sk{प्रतिजानीते}}, \emph{pratijānīte}.) Hunting? (\textbf{\sk{मृगया}}, \emph{mṛgayā}.) Merit? (\textbf{\sk{पुण्यम्}}, \emph{puṇyam}.) A sin? (\textbf{\sk{पापम्}}, \emph{pāpam}.) Heaven? (\textbf{\sk{स्वर्गः}}, \emph{svargaḥ}.) Hell? (\textbf{\sk{नरकः}}, \emph{narakaḥ}.) Reaps? (\textbf{\sk{भुंक्ते}}, \emph{bhuṅkte}.) A crown prince? (\textbf{\sk{युवराजः}}, \emph{yuvarājaḥ}.) A bride's choice? (\textbf{\sk{स्वयंवरः}}, \emph{svayaṃvaraḥ}.) Chooses? (\textbf{\sk{वृणीते}}, \emph{vṛṇīte}.) Disguised? (\textbf{\sk{प्रच्छन्नः}}, \emph{pracchannaḥ}.) Justice? (\textbf{\sk{न्यायः}}, \emph{nyāyaḥ}.)
\end{tcolorbox}
\section[(dialogue)]{The Prince's Choice — once more}

\label{lesson:SA-R339-review-final-339}



\begin{quote}
Say the Sanskrit for a crown prince and a bride's choice.
\end{quote}

\subsection*{Your turn}
\begin{itemize}
\raggedright
  \item \emph{Say it:} a crown prince, a bride's choice, chooses, disguised, justice
\end{itemize}

\begin{tcolorbox}[breakable,colback=teal!4,colframe=teal!35!black,title={Before you move on}]
A crown prince? (\textbf{\sk{युवराजः}}, \emph{yuvarājaḥ}.) A bride's choice? (\textbf{\sk{स्वयंवरः}}, \emph{svayaṃvaraḥ}.) Chooses? (\textbf{\sk{वृणीते}}, \emph{vṛṇīte}.) Disguised? (\textbf{\sk{प्रच्छन्नः}}, \emph{pracchannaḥ}.) Justice? (\textbf{\sk{न्यायः}}, \emph{nyāyaḥ}.)
\end{tcolorbox}
